\section{Conclusion}
\label{sec:conclusion}

We built the composition of sparse expert routing and adaptive recursion as a
single engine, held parameters and inputs fixed across its three configurations,
and measured what the combination buys at each level of inference compute. At
5.6M parameters on WikiText-103 it buys efficiency and robustness rather than
peak quality: all three arms lie on the quality/compute frontier, MoRE buys the
most quality per unit of compute ($7.9\%$ better than MoE at $2.19\times$ its
FLOPs, $72\%$ of recursion's gain for $24\%$ of recursion's compute), and out of
domain it \emph{dominates} recursion alone---lower loss at a quarter of the
compute, the one respect in which composing the mechanisms does something neither
does alone. What it does not buy is in-domain quality: recursion alone remains
lowest-loss by 0.1216 nats, significant at our resolution floor, and routing adds
nothing to routing quality over MoE either. We keep the two axes separate
throughout, and neither statement licenses the other.

The mechanism is not the obvious one. Routing does not dilute the update
budget---tokens per parameter are identical by construction---and the recursion
is not dynamically broken, since subspace trapping fails three ways. What the
evidence supports is a competition for one per-token budget: the arm that
preserves specialisation saturates its depth. We release the engine, the
evaluation battery, and the forced-depth control, and it is the last we would
most like to see reused. A mean step count and an exit histogram do not establish
that depth depends on the input---a policy ignoring its input reproduces
both---whereas scoring against a uniform policy forced to the same mean depth
does, costs one inference pass, and is what separates our two recursive arms by a
factor of 24. We recommend it as a routine check before any model is described as
allocating computation adaptively.
